\section{Threat Model}
\label{sec:threat}

\subsection{System Model and Trust Boundary}
\label{sec:threat:trust}

\SystemName targets LLM-agent settings in which agents invoke tools, shell
commands, scripts, or binaries that access files, spawn processes, use the
network, or interact with external services. The TCB includes the operating
system kernel, the kernel enforcement mechanisms used by \SystemName, the
\SystemName runtime, the SK-IR compiler and rule base, the policy installer, the
ghost-policy manager, ShadowFS, and the semantic validator. We assume these
components are not compromised or attacker-modifiable. Our goal is not to
minimize the TCB, but to exclude the agent framework, generated actions, and
invoked tools from it.

The agent and all code it controls are untrusted, including prompts, model
outputs, tool-selection logic, permission declarations, generated scripts,
plugins, and launched binaries. Host-installed binaries receive no special
trust; they execute under the same invocation-scoped policy as the selected tool.

The semantic validator is trusted to authorize releases using a bounded context
constructed by the runtime, including release-relevant information such as the
user prompt, task status, history operations, and current tool intent. \SystemName
does not infer user intent perfectly: if the trusted validator approves an unsafe
release, \SystemName enforces that decision. The guarantee is that
agent-generated actions cannot exceed the compiled kernel authority or become
externally visible without validator approval.

\subsection{Adversary Capabilities}
\label{sec:threat:adversary}

We consider adversaries that try to make an agent act outside the user's task
semantics or beyond the authority compiled for an invocation. The adversary may
control any agent-observed input, including prompts, documents, webpages,
retrieved content, tool outputs, third-party responses, and data consumed by
launched programs. These inputs may contain prompt injections, malicious
instructions, or malformed data that induce excessive permissions, dangerous
commands, unauthorized file access, sensitive-resource modification, or network
exfiltration.

Any program launched by the agent may be malicious. It may spawn helper
processes, exploit path aliases or symbolic links, access resources through file
descriptors or \texttt{/proc}, exercise ambient privileges, use network APIs for
exfiltration, or otherwise attempt to evade confinement. The adversary may also
craft permission declarations and tool parameters that stress the SK-IR compiler,
including ambiguous paths, broad network destinations, malformed endpoints, and
over-limit requests. \SystemName treats these inputs as untrusted and requires
deterministic validation and fail-closed compilation; the adversary cannot modify
the TCB.

\subsection{Security Goals}
\label{sec:threat:goals}

\SystemName aims to provide the following properties.

\paragraph{Mediation for modeled operations.}
All modeled filesystem, process, privilege, and network operations are mediated
at the operating-system boundary for managed invocations. Requests outside this
modeled surface are rejected during compilation rather than executed with ambient
authority.

\paragraph{Invocation-scoped least privilege.}
Before untrusted tool code runs, \SystemName installs an invocation-specific
kernel policy derived from accepted permissions. The SK-IR compiler may
normalize, validate, or narrow requested authority, but must not expand it;
requests that cannot be safely compiled fail closed. Permissions are bound to
the invocation and revoked on completion, failure, timeout, discard, or lease
expiration.

\paragraph{Fail-closed semantic release.}
Agent-visible outputs, committed filesystem changes, and externally visible
effects are released only after trusted semantic approval. Denial, timeout, or
pre-approval failure discards speculative state and revokes authority.

\paragraph{Controlled irreversible effects.}
Irreversible external effects, such as network transmissions, API calls, and
remote-service interactions, are release boundaries. They must be withheld or
mediated until semantic authorization succeeds.

\paragraph{Low-latency authorization.}
\SystemName overlaps confined execution with semantic validation to reduce
user-visible latency. This performance goal does not relax the above security
properties: speculation changes when work is performed, not when unauthorized
effects become visible.

\subsection{Out of Scope}
\label{sec:threat:out}

\SystemName does not protect against compromise of the operating-system kernel,
trusted kernel mechanisms, runtime, SK-IR compiler, ShadowFS, ghost-policy
manager, or semantic validator. Bugs or compromise in these components may
violate its guarantees. Side channels, hardware attacks, physical attacks,
host-level denial of service, and resource-exhaustion attacks are outside scope;
\SystemName denies release on failure but does not guarantee availability.

\SystemName also does not determine absolute user intent. Incorrect approvals
are trusted policy decisions, while incorrect denials may reduce utility without
violating the stated safety properties. Finally, \SystemName prevents external
effects from being released before authorization, but does not roll back effects
that have already been externally observed.
